\begin{abstract}
\noindent
Large language model (LLM) systems are increasingly proposed to assist scholarly
peer review, yet most evaluations judge the prose of machine-generated review
\emph{text} rather than the validity of the numeric \emph{score} a system assigns.
We validate \textsc{AIPR}, which reads a submitted manuscript and emits five
0--100 quality dimensions (novelty, rigor, applicability, clarity, citation) and a
weighted overall score, against the public decision outcomes of a major machine
learning venue. We validate the overall score; one subscore, citation, is not
validated in this cohort, where an audit fallback pinned it. AIPR is deployed as an interactive assistant; here we validate its
non-interactive baseline score, the first-pass signal produced before any reviewer
engages with it. Across \NminiPrimary{} \mbox{ICLR} submissions with public
decision tiers (reject/poster/oral) and mean reviewer ratings, graded
under a single frozen pipeline version (v6) with hypotheses and thresholds
pre-registered before any score met any outcome, the overall score separates rejected from accepted submissions
(AUROC~\aurocMiniFull), rises monotonically across the three tiers (Spearman
$\rho=\trendRho$, permutation $\trendPrel$), and correlates with the mean reviewer
rating ($\rho =$~\spearmanRatingMiniFull). The signal is strongest where we claim it:
the lowest score quintile is rejected at \bottomRejectRateFrontier\% on the
production model directly (\bottomRejectRate\% at scale on a cheap-model cohort, a
\bottomLift$\times$ lift over the \baseRejectRate\% base rate), with oral papers
absent from that band. Our central finding is about the model, not the pipeline. A
generic one-paragraph prompt on the same frontier model and PDF---no rubric, no
audit---already separates rejected from accepted submissions on its own
(AUROC~\aurocNaiveFull). The full pipeline is higher by
$\Delta$AUROC \aurocDiffFullNaive{} (95\% CI \aurocDiffCI, \aurocDiffPrel), but this
planned comparison is not statistically resolved; the study was not powered to
rule in or rule out such a small difference over an already strong naive baseline.
\emph{Within this ICLR cohort, with or without our prompt, a frontier model already
produces a first-pass reviewing signal that agrees substantially with human
outcomes}: \emph{intelligence is not the observed bottleneck}, and the human keeps
the decision. What the engineered pipeline adds is
therefore not validity but deployability---it matches that score while grading far
more consistently across repeated runs (within-paper SD \fullRunSD{} vs.\
\naiveRunSD{} points) and returns a grounded, rubric-anchored review in a single
pass. We do not claim the score predicts the acceptance decision or replaces reviewers; we
characterize its failure modes directly. The result supports a narrow, defensible
use: an automated first pass that reliably surfaces weak submissions for human
attention.
\end{abstract}
